% ---
% id: foundations/lem:absorbing-domination
% kind: lem
% title: "Domination of a family that absorbs its residue in one element"
% summary: >-
%   When every element but one is rounded down below the reference element and
%   the remaining one takes the rest, only that designated element constrains
%   the weight. A weight below one exists exactly when the absorbed residue lies
%   in its support, and it is at most residue norm over least eigenvalue.
% prominence: standard
% status: accepted
% uses: [def:domination, lem:domination]
% mentions: []
% assumes: []
% refs: []
% provenance: newly-derived
% notes: "No requirement uses this lemma; rem:noise-floor reads its part (iv) for the weight a noise floor allows."
% lean:
% ---

A family that keeps every element but one below its reference element and
lets one designated element take the rest is the simplest family a model can
dominate at a small weight. The lemma shows that the weight of the domination
is then decided by the designated element alone, and how large it is.

\begin{lemma}\label{lem:absorbing-domination}
Let $s$ be a block of dimension $d_s$, $J$ a finite non-empty alphabet,
$c>0$, and $j^{\star}\in J$. Let $\{G_j\}_{j\in J}$ be Hermitian positive
semidefinite $d_s\times d_s$ matrices with $\sum_{j\in J}G_j=c\,\idon{d_s}$,
and let $\{F_j\}_{j\in J}$ be Hermitian matrices with
\begin{equation}\label{eq:absorbing}
  0\preceq F_j\preceq G_j\quad(j\neq j^{\star}),
  \qquad
  F_{j^{\star}}=c\,\idon{d_s}-\sum_{j\neq j^{\star}}F_j .
\end{equation}
Put $\Delta\coloneqq\sum_{j\neq j^{\star}}(G_j-F_j)$, the \emph{absorbed
residue}.
\begin{enumerate}[label=(\roman*),leftmargin=2.4em]
\item $\Delta\succeq0$, $F_{j^{\star}}=G_{j^{\star}}+\Delta\succeq G_{j^{\star}}$,
  every $F_j$ is positive semidefinite, and $\sum_jF_j=c\,\idon{d_s}$.
\item For $q\in[0,1)$, $\{G_j\}$ dominates $\{F_j\}$ at $q$ on this block
  (Definition~\ref{def:domination}) if and only if
  \begin{equation}\label{eq:absorbing-condition}
    q\,G_{j^{\star}}-(1-q)\,\Delta\;\succeq\;0 .
  \end{equation}
  For every $j\neq j^{\star}$ and every $q\in[0,1)$,
  $G_j-(1-q)F_j\succeq0$ holds.
\item Some $q\in[0,1)$ satisfies \eqref{eq:absorbing-condition} if and only if
  $\mathrm{supp}\,\Delta\subseteq\mathrm{supp}\,G_{j^{\star}}$; this holds in
  particular when $G_{j^{\star}}$ is positive definite.
\item Suppose $\mathrm{supp}\,\Delta\subseteq\mathrm{supp}\,G_{j^{\star}}$,
  let $\lambda>0$ be the least non-zero eigenvalue of $G_{j^{\star}}$ (its
  least eigenvalue when it is positive definite) and let
  $\delta\ge\opn{\Delta}$. If $G_{j^{\star}}\neq0$, then
  \eqref{eq:absorbing-condition} holds for every $q$ with
  $\delta/(\lambda+\delta)\le q<1$. If $\Delta=0$, it holds at $q=0$.
\end{enumerate}
\end{lemma}

\begin{proof}
(i) Each $G_j-F_j$ with $j\neq j^{\star}$ is positive semidefinite by
\eqref{eq:absorbing}, so $\Delta\succeq0$. By completeness of $\{G_j\}$,
$G_{j^{\star}}=c\,\idon{d_s}-\sum_{j\neq j^{\star}}G_j$, and subtracting this
from the definition of $F_{j^{\star}}$ gives
$F_{j^{\star}}-G_{j^{\star}}=\sum_{j\neq j^{\star}}(G_j-F_j)=\Delta$. So
$F_{j^{\star}}=G_{j^{\star}}+\Delta\succeq G_{j^{\star}}\succeq0$; the other
$F_j$ are positive semidefinite by \eqref{eq:absorbing}, and
$\sum_jF_j=c\,\idon{d_s}$ by the definition of $F_{j^{\star}}$.

(ii) For $j\neq j^{\star}$,
$G_j-(1-q)F_j=(G_j-F_j)+qF_j$ is a sum of two positive semidefinite
matrices, by \eqref{eq:absorbing} and $q\ge0$. For $j^{\star}$, by (i),
$G_{j^{\star}}-(1-q)F_{j^{\star}}=G_{j^{\star}}-(1-q)(G_{j^{\star}}+\Delta)
=qG_{j^{\star}}-(1-q)\Delta$. So \eqref{eq:domination} holds at every outcome
if and only if it holds at $j^{\star}$, which is
\eqref{eq:absorbing-condition}.

(iii) If \eqref{eq:absorbing-condition} holds for some $q<1$, then, by (ii),
$\{G_j\}$ dominates $\{F_j\}$ at $q$, and
Lemma~\ref{lem:domination}\,(ii) gives
$\mathrm{supp}\,F_{j^{\star}}\subseteq\mathrm{supp}\,G_{j^{\star}}$. Since
$0\preceq\Delta\preceq\Delta+G_{j^{\star}}=F_{j^{\star}}$, every
$x\in\ker F_{j^{\star}}$ satisfies $0\le\langle x,\Delta x\rangle\le
\langle x,F_{j^{\star}}x\rangle=0$, so $\Delta x=0$ as in the proof of
Lemma~\ref{lem:domination}\,(ii); hence
$\mathrm{supp}\,\Delta\subseteq\mathrm{supp}\,F_{j^{\star}}\subseteq
\mathrm{supp}\,G_{j^{\star}}$. The converse is (iv) when $G_{j^{\star}}\neq0$
and $q=0$ when $\Delta=0$; if $G_{j^{\star}}=0$, the inclusion forces
$\Delta=0$. A positive definite $G_{j^{\star}}$ has support the whole space.

(iv) If $\Delta=0$, \eqref{eq:absorbing-condition} at $q=0$ reads $0\succeq0$.
Now let $G_{j^{\star}}\neq0$, and let $\Pi$ be the orthogonal projector onto
$\mathrm{supp}\,G_{j^{\star}}$. By the spectral decomposition
$G_{j^{\star}}\succeq\lambda\Pi$. Since $\mathrm{supp}\,\Delta\subseteq
\mathrm{supp}\,G_{j^{\star}}$ and $\Delta$ is Hermitian,
$\Delta=\Pi\Delta\Pi\preceq\opn{\Delta}\,\Pi\preceq\delta\,\Pi$. For
$q\ge\delta/(\lambda+\delta)$ we have $q\lambda\ge(1-q)\delta$, so
\begin{equation*}
  qG_{j^{\star}}-(1-q)\Delta\;\succeq\;q\lambda\,\Pi-(1-q)\delta\,\Pi
  \;=\;\big(q\lambda-(1-q)\delta\big)\Pi\;\succeq\;0 .
\end{equation*}
Since $\lambda>0$, $\delta/(\lambda+\delta)<1$, so such $q$ exist in
$[0,1)$.
\end{proof}
